% GENERATED FILE. Edit canonical lessons, then run npm run generate:books.
% canonical-source-hash: fnv1a64:f8ddf46557127118
% canonical-lessons: AR-C87-ams, AR-C87-ghadan, AR-C87-alan, AR-C87-ahyanan, AR-C87-qabla

\chapter{Yesterday, Tomorrow, Now, Sometimes, Before}
\label{ch:ar-yesterday-tomorrow-now-sometimes-before}
\hlchaptermodality{87}

\begin{chapteropening}
\textbf{By the end of this chapter:} \emph{I can say yesterday, tomorrow, now, sometimes and before.}

Complete the last lesson of chapter 87: I can say yesterday, tomorrow, now, sometimes and before.
\end{chapteropening}

\section[\texorpdfstring{\ar{أمس} (ams)}{أمس (ams)}]{\ar{أمس} (ams) — yesterday}

\label{lesson:AR-C87-ams}



\begin{quote}
Before the new one: say the Arabic for a month, then the Arabic for a minute.
\end{quote}

\subsection*{You'll want to know: \ar{أمس}}
\textbf{\ar{أمس}} — \emph{ams} — \textquotedblleft{}yesterday\textquotedblright{}.

\begin{cousinweb}
\textbf{\ar{أمس}} (\emph{ams}) is built on the native Arabic root \emph{ʾ-m-s}, \textquotedblleft{}yesterday\textquotedblright{} (Lane). Hebrew \emph{emesh}, \textquotedblleft{}last night\textquotedblright{}, is usually compared.
\end{cousinweb}

\begin{grammarlens}[title={What you've built}]
One more word for this chapter.
\end{grammarlens}

\subsection*{Your turn}
\begin{itemize}
\raggedright
  \item \emph{Say it:} \emph{ams}
  \item \emph{Say it:} \emph{ams}, once more
  \item \emph{Say it:} say \emph{ams}
  \item \emph{Recall:} say the Arabic for a month, then the Arabic for a minute, then say \emph{ams} again
\end{itemize}

\begin{tcolorbox}[breakable,colback=teal!4,colframe=teal!35!black,title={Before you move on}]
What does \ar{أمس} mean? (\textquotedblleft{}Yesterday\textquotedblright{}.) Say it once more.
\end{tcolorbox}
\section[\texorpdfstring{\ar{غدا} (ghadan)}{غدا (ghadan)}]{\ar{غدا} (ghadan) — tomorrow}

\label{lesson:AR-C87-ghadan}



\begin{quote}
Before the new one: say the Arabic for a minute, then the Arabic for yesterday.
\end{quote}

\subsection*{You'll want to know: \ar{غدا}}
\textbf{\ar{غدا}} — \emph{ghadan} — \textquotedblleft{}tomorrow\textquotedblright{}.

\begin{cousinweb}
\textbf{\ar{غدا}} (\emph{ghadan}) is built on the native Arabic root \emph{gh-d-w}, \textquotedblleft{}to go forth early in the morning\textquotedblright{} (Lane). Lane gives \emph{ghad}, \textquotedblleft{}tomorrow\textquotedblright{}, and \emph{ghudwa}, \textquotedblleft{}early morning\textquotedblright{}, on one root.
\end{cousinweb}

\begin{grammarlens}[title={What you've built}]
One more word for this chapter.
\end{grammarlens}

\subsection*{Your turn}
\begin{itemize}
\raggedright
  \item \emph{Say it:} \emph{ghadan}
  \item \emph{Say it:} \emph{ghadan}, once more
  \item \emph{Say it:} say \emph{ghadan}
  \item \emph{Recall:} say the Arabic for a minute, then the Arabic for yesterday, then say \emph{ghadan} again
\end{itemize}

\begin{tcolorbox}[breakable,colback=teal!4,colframe=teal!35!black,title={Before you move on}]
What does \ar{غدا} mean? (\textquotedblleft{}Tomorrow\textquotedblright{}.) Say it once more.
\end{tcolorbox}
\section[\texorpdfstring{\ar{الآن} (al-ān)}{الآن (al-ān)}]{\ar{الآن} (al-ān) — now}

\label{lesson:AR-C87-alan}



\begin{quote}
Before the new one: say the Arabic for yesterday, then the Arabic for tomorrow.
\end{quote}

\subsection*{You'll want to know: \ar{الآن}}
\textbf{\ar{الآن}} — \emph{al-ān} — \textquotedblleft{}now\textquotedblright{}.

\begin{cousinweb}
\textbf{\ar{الآن}} (\emph{al-ān}) is listed by Lane under \emph{ʾ-y-n} (some put it under \emph{ʾ-w-n}) as a noun of the present time whose \emph{al-} is never removed. It is built on \emph{ān}, \textquotedblleft{}time\textquotedblright{}.
\end{cousinweb}

\begin{grammarlens}[title={What you've built}]
One more word for this chapter.
\end{grammarlens}

\subsection*{Your turn}
\begin{itemize}
\raggedright
  \item \emph{Say it:} \emph{al-ān}
  \item \emph{Say it:} \emph{al-ān}, once more
  \item \emph{Say it:} say \emph{al-ān}
  \item \emph{Recall:} say the Arabic for yesterday, then the Arabic for tomorrow, then say \emph{al-ān} again
\end{itemize}

\begin{tcolorbox}[breakable,colback=teal!4,colframe=teal!35!black,title={Before you move on}]
What does \ar{الآن} mean? (\textquotedblleft{}Now\textquotedblright{}.) Say it once more.
\end{tcolorbox}
\section[\texorpdfstring{\ar{أحيانا} (aḥyānan)}{أحيانا (aḥyānan)}]{\ar{أحيانا} (aḥyānan) — sometimes}

\label{lesson:AR-C87-ahyanan}



\begin{quote}
Before the new one: say the Arabic for tomorrow, then the Arabic for now.
\end{quote}

\subsection*{You'll want to know: \ar{أحيانا}}
\textbf{\ar{أحيانا}} — \emph{aḥyānan} — \textquotedblleft{}sometimes\textquotedblright{}.

\begin{cousinweb}
\textbf{\ar{أحيانا}} (\emph{aḥyānan}) is built on the native Arabic root \emph{ḥ-y-n}, \textquotedblleft{}a time\textquotedblright{} (Lane). \emph{Aḥyānan} is \emph{aḥyān}, the plural of \emph{ḥīn}, \textquotedblleft{}a time\textquotedblright{}, with the adverb ending \emph{-an}: \textquotedblleft{}at times\textquotedblright{}.
\end{cousinweb}

\begin{grammarlens}[title={What you've built}]
One more word for this chapter.
\end{grammarlens}

\subsection*{Your turn}
\begin{itemize}
\raggedright
  \item \emph{Say it:} \emph{aḥyānan}
  \item \emph{Say it:} \emph{aḥyānan}, once more
  \item \emph{Say it:} say \emph{aḥyānan}
  \item \emph{Recall:} say the Arabic for tomorrow, then the Arabic for now, then say \emph{aḥyānan} again
\end{itemize}

\begin{tcolorbox}[breakable,colback=teal!4,colframe=teal!35!black,title={Before you move on}]
What does \ar{أحيانا} mean? (\textquotedblleft{}Sometimes\textquotedblright{}.) Say it once more.
\end{tcolorbox}
\section[\texorpdfstring{\ar{قبل} (qabla)}{قبل (qabla)}]{\ar{قبل} (qabla) — before}

\label{lesson:AR-C87-qabla}



\begin{quote}
Before the new one: say the Arabic for now, then the Arabic for sometimes.
\end{quote}

\subsection*{You'll want to know: \ar{قبل}}
\textbf{\ar{قبل}} — \emph{qabla} — \textquotedblleft{}before\textquotedblright{}.

\begin{cousinweb}
\textbf{\ar{قبل}} (\emph{qabla}) is built on the native Arabic root \emph{q-b-l}, \textquotedblleft{}to face, to be in front\textquotedblright{} (Lane). The same root gives \emph{qibla}, the direction one faces in prayer.
\end{cousinweb}

\begin{grammarlens}[title={What you've built}]
One more word for this chapter.
\end{grammarlens}

\subsection*{Your turn}
\begin{itemize}
\raggedright
  \item \emph{Say it:} \emph{qabla}
  \item \emph{Say it:} \emph{qabla}, once more
  \item \emph{Say it:} say \emph{qabla}
  \item \emph{Recall:} say the Arabic for now, then the Arabic for sometimes, then say \emph{qabla} again
\end{itemize}

\begin{tcolorbox}[breakable,colback=teal!4,colframe=teal!35!black,title={Before you move on}]
What does \ar{قبل} mean? (\textquotedblleft{}Before\textquotedblright{}.) Say it once more.
\end{tcolorbox}
